% Estilo comun de los manuales, el mismo de la propuesta y del reporte de avances:
% fondo oscuro, titulos morados y recuadro de encabezado. Se carga con % Estilo comun de los manuales, el mismo de la propuesta y del reporte de avances:
% fondo oscuro, titulos morados y recuadro de encabezado. Se carga con % Estilo comun de los manuales, el mismo de la propuesta y del reporte de avances:
% fondo oscuro, titulos morados y recuadro de encabezado. Se carga con % Estilo comun de los manuales, el mismo de la propuesta y del reporte de avances:
% fondo oscuro, titulos morados y recuadro de encabezado. Se carga con \input{estilo}.
\usepackage[utf8]{inputenc}
\usepackage[T1]{fontenc}
\usepackage[spanish,es-noshorthands,es-tabla]{babel}
\usepackage{lmodern}
\usepackage[margin=2.5cm]{geometry}
\usepackage{graphicx}
\usepackage{subcaption}
\usepackage{float}
\usepackage[table]{xcolor}
\usepackage{titlesec}
\usepackage{tcolorbox}
\usepackage{enumitem}
\usepackage{array}
\usepackage{longtable}
\usepackage{amssymb}
\usepackage{pgfgantt}
\usepackage{hyperref}

\definecolor{fondo}{HTML}{1E1E2E}
\definecolor{texto}{HTML}{CDCDD6}
\definecolor{morado}{HTML}{B48EDB}
\definecolor{azul}{HTML}{8AADF4}
\definecolor{rosa}{HTML}{E3A8D6}
\definecolor{verde}{HTML}{A6DA95}
\definecolor{linea}{HTML}{55556A}

\pagecolor{fondo}
\color{texto}
\hypersetup{colorlinks=true, linkcolor=azul, urlcolor=azul, citecolor=azul}
\arrayrulecolor{texto}
\graphicspath{{../docs/evidencias/}{img/}}

\setlength{\parindent}{0pt}
\setlength{\parskip}{6pt}

% secciones numeradas, moradas y con una linea abajo
\titleformat{\section}{\Large\bfseries\color{morado}}{\thesection.}{8pt}{}[{\color{linea}\titlerule}]
\titleformat{\subsection}{\large\bfseries\color{azul}}{\thesubsection}{8pt}{}
\titleformat{\subsubsection}{\normalsize\bfseries\color{rosa}}{}{0pt}{}
\titlespacing*{\section}{0pt}{18pt}{10pt}
\titlespacing*{\subsection}{0pt}{12pt}{6pt}

\setlist[itemize,1]{label=\scriptsize$\blacksquare$}
\setlist{itemsep=3pt}

% lo que falta por escribir o por confirmar; se busca con "pendiente" para no dejar ninguno
\newcommand{\pendiente}[1]{%
  \begin{tcolorbox}[colback=fondo, colframe=rosa, coltext=rosa, arc=2mm, boxrule=0.8pt, left=6pt, right=6pt, top=4pt, bottom=4pt]
    \textbf{Pendiente:} #1
  \end{tcolorbox}}

% tecla con recuadro
\newcommand{\tecla}[1]{\tcbox[on line, colback=fondo, colframe=azul, coltext=azul, arc=1mm, boxrule=0.6pt, left=2pt, right=2pt, top=0pt, bottom=0pt]{\small\ttfamily #1}}

\newcommand{\dosfotos}[5]{%
  \begin{figure}[H]
    \centering
    \begin{subfigure}{0.48\textwidth}
      \includegraphics[width=\linewidth]{#1}
      \caption{#2}
    \end{subfigure}\hfill
    \begin{subfigure}{0.48\textwidth}
      \includegraphics[width=\linewidth]{#3}
      \caption{#4}
    \end{subfigure}
    \caption{#5}
  \end{figure}}

\newcommand{\unafoto}[2]{%
  \begin{figure}[H]
    \centering
    \includegraphics[width=0.8\textwidth]{#1}
    \caption{#2}
  \end{figure}}

% encabezado en recuadro: titulo del documento
\newcommand{\encabezado}[1]{%
  \begin{tcolorbox}[colback=fondo, colframe=morado, coltext=texto, arc=4mm, boxrule=0.8pt, left=10pt, right=10pt, top=8pt, bottom=8pt]
    \centering
    {\Large\bfseries\color{azul} #1 $\|$ Calenda de Oaxaca}\\[4pt]
    {\large\bfseries Facultad de Ingeniería -- UNAM}\\[10pt]
    \begin{tabular}{@{}p{0.47\linewidth}p{0.47\linewidth}@{}}
      {\color{rosa}\textbf{Equipo:}} 16 & {\color{rosa}\textbf{Materia:}} Computación Gráfica e I.H.C. \\
      {\color{rosa}\textbf{Número de cuenta:}} 421112601 & {\color{rosa}\textbf{Grupo:}} 01 \\
      {\color{rosa}\textbf{Semestre:}} 2027-1 & {\color{rosa}\textbf{Entrega:}} 18 de noviembre de 2026 \\
    \end{tabular}
  \end{tcolorbox}}
.
\usepackage[utf8]{inputenc}
\usepackage[T1]{fontenc}
\usepackage[spanish,es-noshorthands,es-tabla]{babel}
\usepackage{lmodern}
\usepackage[margin=2.5cm]{geometry}
\usepackage{graphicx}
\usepackage{subcaption}
\usepackage{float}
\usepackage[table]{xcolor}
\usepackage{titlesec}
\usepackage{tcolorbox}
\usepackage{enumitem}
\usepackage{array}
\usepackage{longtable}
\usepackage{amssymb}
\usepackage{pgfgantt}
\usepackage{hyperref}

\definecolor{fondo}{HTML}{1E1E2E}
\definecolor{texto}{HTML}{CDCDD6}
\definecolor{morado}{HTML}{B48EDB}
\definecolor{azul}{HTML}{8AADF4}
\definecolor{rosa}{HTML}{E3A8D6}
\definecolor{verde}{HTML}{A6DA95}
\definecolor{linea}{HTML}{55556A}

\pagecolor{fondo}
\color{texto}
\hypersetup{colorlinks=true, linkcolor=azul, urlcolor=azul, citecolor=azul}
\arrayrulecolor{texto}
\graphicspath{{../docs/evidencias/}{img/}}

\setlength{\parindent}{0pt}
\setlength{\parskip}{6pt}

% secciones numeradas, moradas y con una linea abajo
\titleformat{\section}{\Large\bfseries\color{morado}}{\thesection.}{8pt}{}[{\color{linea}\titlerule}]
\titleformat{\subsection}{\large\bfseries\color{azul}}{\thesubsection}{8pt}{}
\titleformat{\subsubsection}{\normalsize\bfseries\color{rosa}}{}{0pt}{}
\titlespacing*{\section}{0pt}{18pt}{10pt}
\titlespacing*{\subsection}{0pt}{12pt}{6pt}

\setlist[itemize,1]{label=\scriptsize$\blacksquare$}
\setlist{itemsep=3pt}

% lo que falta por escribir o por confirmar; se busca con "pendiente" para no dejar ninguno
\newcommand{\pendiente}[1]{%
  \begin{tcolorbox}[colback=fondo, colframe=rosa, coltext=rosa, arc=2mm, boxrule=0.8pt, left=6pt, right=6pt, top=4pt, bottom=4pt]
    \textbf{Pendiente:} #1
  \end{tcolorbox}}

% tecla con recuadro
\newcommand{\tecla}[1]{\tcbox[on line, colback=fondo, colframe=azul, coltext=azul, arc=1mm, boxrule=0.6pt, left=2pt, right=2pt, top=0pt, bottom=0pt]{\small\ttfamily #1}}

\newcommand{\dosfotos}[5]{%
  \begin{figure}[H]
    \centering
    \begin{subfigure}{0.48\textwidth}
      \includegraphics[width=\linewidth]{#1}
      \caption{#2}
    \end{subfigure}\hfill
    \begin{subfigure}{0.48\textwidth}
      \includegraphics[width=\linewidth]{#3}
      \caption{#4}
    \end{subfigure}
    \caption{#5}
  \end{figure}}

\newcommand{\unafoto}[2]{%
  \begin{figure}[H]
    \centering
    \includegraphics[width=0.8\textwidth]{#1}
    \caption{#2}
  \end{figure}}

% encabezado en recuadro: titulo del documento
\newcommand{\encabezado}[1]{%
  \begin{tcolorbox}[colback=fondo, colframe=morado, coltext=texto, arc=4mm, boxrule=0.8pt, left=10pt, right=10pt, top=8pt, bottom=8pt]
    \centering
    {\Large\bfseries\color{azul} #1 $\|$ Calenda de Oaxaca}\\[4pt]
    {\large\bfseries Facultad de Ingeniería -- UNAM}\\[10pt]
    \begin{tabular}{@{}p{0.47\linewidth}p{0.47\linewidth}@{}}
      {\color{rosa}\textbf{Equipo:}} 16 & {\color{rosa}\textbf{Materia:}} Computación Gráfica e I.H.C. \\
      {\color{rosa}\textbf{Número de cuenta:}} 421112601 & {\color{rosa}\textbf{Grupo:}} 01 \\
      {\color{rosa}\textbf{Semestre:}} 2027-1 & {\color{rosa}\textbf{Entrega:}} 18 de noviembre de 2026 \\
    \end{tabular}
  \end{tcolorbox}}
.
\usepackage[utf8]{inputenc}
\usepackage[T1]{fontenc}
\usepackage[spanish,es-noshorthands,es-tabla]{babel}
\usepackage{lmodern}
\usepackage[margin=2.5cm]{geometry}
\usepackage{graphicx}
\usepackage{subcaption}
\usepackage{float}
\usepackage[table]{xcolor}
\usepackage{titlesec}
\usepackage{tcolorbox}
\usepackage{enumitem}
\usepackage{array}
\usepackage{longtable}
\usepackage{amssymb}
\usepackage{pgfgantt}
\usepackage{hyperref}

\definecolor{fondo}{HTML}{1E1E2E}
\definecolor{texto}{HTML}{CDCDD6}
\definecolor{morado}{HTML}{B48EDB}
\definecolor{azul}{HTML}{8AADF4}
\definecolor{rosa}{HTML}{E3A8D6}
\definecolor{verde}{HTML}{A6DA95}
\definecolor{linea}{HTML}{55556A}

\pagecolor{fondo}
\color{texto}
\hypersetup{colorlinks=true, linkcolor=azul, urlcolor=azul, citecolor=azul}
\arrayrulecolor{texto}
\graphicspath{{../docs/evidencias/}{img/}}

\setlength{\parindent}{0pt}
\setlength{\parskip}{6pt}

% secciones numeradas, moradas y con una linea abajo
\titleformat{\section}{\Large\bfseries\color{morado}}{\thesection.}{8pt}{}[{\color{linea}\titlerule}]
\titleformat{\subsection}{\large\bfseries\color{azul}}{\thesubsection}{8pt}{}
\titleformat{\subsubsection}{\normalsize\bfseries\color{rosa}}{}{0pt}{}
\titlespacing*{\section}{0pt}{18pt}{10pt}
\titlespacing*{\subsection}{0pt}{12pt}{6pt}

\setlist[itemize,1]{label=\scriptsize$\blacksquare$}
\setlist{itemsep=3pt}

% lo que falta por escribir o por confirmar; se busca con "pendiente" para no dejar ninguno
\newcommand{\pendiente}[1]{%
  \begin{tcolorbox}[colback=fondo, colframe=rosa, coltext=rosa, arc=2mm, boxrule=0.8pt, left=6pt, right=6pt, top=4pt, bottom=4pt]
    \textbf{Pendiente:} #1
  \end{tcolorbox}}

% tecla con recuadro
\newcommand{\tecla}[1]{\tcbox[on line, colback=fondo, colframe=azul, coltext=azul, arc=1mm, boxrule=0.6pt, left=2pt, right=2pt, top=0pt, bottom=0pt]{\small\ttfamily #1}}

\newcommand{\dosfotos}[5]{%
  \begin{figure}[H]
    \centering
    \begin{subfigure}{0.48\textwidth}
      \includegraphics[width=\linewidth]{#1}
      \caption{#2}
    \end{subfigure}\hfill
    \begin{subfigure}{0.48\textwidth}
      \includegraphics[width=\linewidth]{#3}
      \caption{#4}
    \end{subfigure}
    \caption{#5}
  \end{figure}}

\newcommand{\unafoto}[2]{%
  \begin{figure}[H]
    \centering
    \includegraphics[width=0.8\textwidth]{#1}
    \caption{#2}
  \end{figure}}

% encabezado en recuadro: titulo del documento
\newcommand{\encabezado}[1]{%
  \begin{tcolorbox}[colback=fondo, colframe=morado, coltext=texto, arc=4mm, boxrule=0.8pt, left=10pt, right=10pt, top=8pt, bottom=8pt]
    \centering
    {\Large\bfseries\color{azul} #1 $\|$ Calenda de Oaxaca}\\[4pt]
    {\large\bfseries Facultad de Ingeniería -- UNAM}\\[10pt]
    \begin{tabular}{@{}p{0.47\linewidth}p{0.47\linewidth}@{}}
      {\color{rosa}\textbf{Equipo:}} 16 & {\color{rosa}\textbf{Materia:}} Computación Gráfica e I.H.C. \\
      {\color{rosa}\textbf{Número de cuenta:}} 421112601 & {\color{rosa}\textbf{Grupo:}} 01 \\
      {\color{rosa}\textbf{Semestre:}} 2027-1 & {\color{rosa}\textbf{Entrega:}} 18 de noviembre de 2026 \\
    \end{tabular}
  \end{tcolorbox}}
.
\usepackage[utf8]{inputenc}
\usepackage[T1]{fontenc}
\usepackage[spanish,es-noshorthands,es-tabla]{babel}
\usepackage{lmodern}
\usepackage[margin=2.5cm]{geometry}
\usepackage{graphicx}
\usepackage{subcaption}
\usepackage{float}
\usepackage[table]{xcolor}
\usepackage{titlesec}
\usepackage{tcolorbox}
\usepackage{enumitem}
\usepackage{array}
\usepackage{longtable}
\usepackage{amssymb}
\usepackage{pgfgantt}
\usepackage{hyperref}

\definecolor{fondo}{HTML}{1E1E2E}
\definecolor{texto}{HTML}{CDCDD6}
\definecolor{morado}{HTML}{B48EDB}
\definecolor{azul}{HTML}{8AADF4}
\definecolor{rosa}{HTML}{E3A8D6}
\definecolor{verde}{HTML}{A6DA95}
\definecolor{linea}{HTML}{55556A}

\pagecolor{fondo}
\color{texto}
\hypersetup{colorlinks=true, linkcolor=azul, urlcolor=azul, citecolor=azul}
\arrayrulecolor{texto}
\graphicspath{{../docs/evidencias/}{img/}}

\setlength{\parindent}{0pt}
\setlength{\parskip}{6pt}

% secciones numeradas, moradas y con una linea abajo
\titleformat{\section}{\Large\bfseries\color{morado}}{\thesection.}{8pt}{}[{\color{linea}\titlerule}]
\titleformat{\subsection}{\large\bfseries\color{azul}}{\thesubsection}{8pt}{}
\titleformat{\subsubsection}{\normalsize\bfseries\color{rosa}}{}{0pt}{}
\titlespacing*{\section}{0pt}{18pt}{10pt}
\titlespacing*{\subsection}{0pt}{12pt}{6pt}

\setlist[itemize,1]{label=\scriptsize$\blacksquare$}
\setlist{itemsep=3pt}

% lo que falta por escribir o por confirmar; se busca con "pendiente" para no dejar ninguno
\newcommand{\pendiente}[1]{%
  \begin{tcolorbox}[colback=fondo, colframe=rosa, coltext=rosa, arc=2mm, boxrule=0.8pt, left=6pt, right=6pt, top=4pt, bottom=4pt]
    \textbf{Pendiente:} #1
  \end{tcolorbox}}

% tecla con recuadro
\newcommand{\tecla}[1]{\tcbox[on line, colback=fondo, colframe=azul, coltext=azul, arc=1mm, boxrule=0.6pt, left=2pt, right=2pt, top=0pt, bottom=0pt]{\small\ttfamily #1}}

\newcommand{\dosfotos}[5]{%
  \begin{figure}[H]
    \centering
    \begin{subfigure}{0.48\textwidth}
      \includegraphics[width=\linewidth]{#1}
      \caption{#2}
    \end{subfigure}\hfill
    \begin{subfigure}{0.48\textwidth}
      \includegraphics[width=\linewidth]{#3}
      \caption{#4}
    \end{subfigure}
    \caption{#5}
  \end{figure}}

\newcommand{\unafoto}[2]{%
  \begin{figure}[H]
    \centering
    \includegraphics[width=0.8\textwidth]{#1}
    \caption{#2}
  \end{figure}}

% encabezado en recuadro: titulo del documento
\newcommand{\encabezado}[1]{%
  \begin{tcolorbox}[colback=fondo, colframe=morado, coltext=texto, arc=4mm, boxrule=0.8pt, left=10pt, right=10pt, top=8pt, bottom=8pt]
    \centering
    {\Large\bfseries\color{azul} #1 $\|$ Calenda de Oaxaca}\\[4pt]
    {\large\bfseries Facultad de Ingeniería -- UNAM}\\[10pt]
    \begin{tabular}{@{}p{0.47\linewidth}p{0.47\linewidth}@{}}
      {\color{rosa}\textbf{Equipo:}} 16 & {\color{rosa}\textbf{Materia:}} Computación Gráfica e I.H.C. \\
      {\color{rosa}\textbf{Número de cuenta:}} 421112601 & {\color{rosa}\textbf{Grupo:}} 01 \\
      {\color{rosa}\textbf{Semestre:}} 2027-1 & {\color{rosa}\textbf{Entrega:}} 18 de noviembre de 2026 \\
    \end{tabular}
  \end{tcolorbox}}
